\subsection{Обоснование метода потенциалов}
\label{subsec:potentials-theory}

\begin{idea}[title={Откуда этот материал}]
Теорема об оптимальности плана и оценка улучшения на лекции не формулировались
— разбирался только сам алгоритм (§7). Ниже обе недостающие части
восстановлены; они закрывают экзаменационный вопрос о методе потенциалов.
\end{idea}

\begin{theorem}[об оптимальности плана]\label{thm:potentials-opt}
Пусть для допустимого базисного плана $X = (x_{ij})$ найдены потенциалы
$u_i$ и $v_j$ такие, что
\[
  v_j - u_i = c_{ij} \ \text{для заполненных клеток},
  \qquad
  v_j - u_i \le c_{ij} \ \text{для пустых} .
\]
Тогда план $X$ оптимален.
\end{theorem}

\begin{proof}
Рассмотрим произвольный допустимый план $Y = (y_{ij})$:
\[
  \sum_{j=1}^{n} y_{ij} = a_i , \qquad
  \sum_{i=1}^{m} y_{ij} = b_j , \qquad y_{ij} \ge 0 .
\]
Для всех клеток выполнено $v_j - u_i \le c_{ij}$, то есть
$c_{ij} \ge v_j - u_i$. Умножая на $y_{ij} \ge 0$ и суммируя,
\[
  \sum_{i=1}^{m}\sum_{j=1}^{n} c_{ij} y_{ij}
  \ \ge\ \sum_{i=1}^{m}\sum_{j=1}^{n} (v_j - u_i)\, y_{ij}
  = \sum_{j=1}^{n} v_j \underbrace{\sum_{i=1}^{m} y_{ij}}_{=\ b_j}
    - \sum_{i=1}^{m} u_i \underbrace{\sum_{j=1}^{n} y_{ij}}_{=\ a_i} ,
\]
откуда
\begin{equation}\label{eq:pot-bound}
  F(Y) \ \ge\ \sum_{j=1}^{n} v_j b_j - \sum_{i=1}^{m} u_i a_i .
\end{equation}
Теперь возьмём сам план $X$. Для заполненных клеток $v_j - u_i = c_{ij}$, а для
пустых $x_{ij} = 0$, поэтому в сумме можно заменить $c_{ij}$ на $v_j - u_i$
\emph{во всех} клетках:
\[
  F(X) = \sum_{i=1}^{m}\sum_{j=1}^{n} c_{ij} x_{ij}
       = \sum_{i=1}^{m}\sum_{j=1}^{n} (v_j - u_i)\, x_{ij}
       = \sum_{j=1}^{n} v_j b_j - \sum_{i=1}^{m} u_i a_i .
\]
Сравнивая с~\eqref{eq:pot-bound}, получаем $F(Y) \ge F(X)$ для любого
допустимого $Y$, то есть $X$ оптимален.
\end{proof}

\begin{idea}
Величина $\sum_j v_j b_j - \sum_i u_i a_i$ — значение двойственной задачи
(см.~\S\ref{subsec:transport-duality-anchor}); неравенство~\eqref{eq:pot-bound}
есть слабая двойственность, а совпадение на плане $X$ — сильная. Поэтому
проверка $\Delta_{ij} = v_j - u_i - c_{ij} \le 0$ и есть проверка двойственной
допустимости.
\end{idea}

\subsubsection*{Насколько улучшается план при пересчёте}

\begin{theorem}\label{thm:potentials-improve}
Пусть в базис вводится пустая клетка $(p,q)$, для которой
$\Delta_{pq} = v_q - u_p - c_{pq} > 0$, и построен цикл пересчёта с величиной
шага $\theta > 0$. Тогда изменение целевой функции равно
\[
  \Delta F = -\,\theta\, \Delta_{pq} \ <\ 0 .
\]
\end{theorem}

\begin{proof}
Обозначим через $C_+$ и $C_-$ множества клеток цикла со знаками $\oplus$ и
$\ominus$. К клеткам из $C_+$ прибавляется $\theta$, из $C_-$ вычитается,
поэтому
\[
  \Delta F = \theta \Big( \sum_{(i,j) \in C_+} c_{ij}
                        - \sum_{(i,j) \in C_-} c_{ij} \Big) .
\]
Все клетки цикла, кроме стартовой $(p,q)$, заполнены, значит для них
$c_{ij} = v_j - u_i$; для $(p,q)$ имеем
$c_{pq} = v_q - u_p - \Delta_{pq}$. Подставим и заметим, что цикл замкнут: в
нём каждая задействованная строка и каждый задействованный столбец встречаются
ровно дважды — один раз со знаком $\oplus$ и один раз со знаком $\ominus$,
поэтому все $u_i$ и $v_j$ взаимно сокращаются. Остаётся только добавка от
стартовой клетки:
\[
  \sum_{(i,j)\in C_+} c_{ij} - \sum_{(i,j)\in C_-} c_{ij} = -\,\Delta_{pq} ,
  \qquad\text{откуда}\qquad
  \Delta F = -\,\theta\,\Delta_{pq} .
\]
Так как $\theta > 0$ и $\Delta_{pq} > 0$, получаем $\Delta F < 0$.
\end{proof}

\begin{idea}
Эта формула и была проверена численно в примере §7: там выбиралась клетка с
$\Delta_{14} = -13$, и стоимость выросла ровно на
$-\theta\Delta_{14} = 13$; а правильная клетка $(3,1)$ с $\Delta_{31} = +6$ и
$\theta = 6$ дала падение на $36$. Отсюда же правило выбора: брать клетку с
\emph{наибольшим} $\Delta_{ij} > 0$ — при прочих равных это даёт наибольшее
уменьшение стоимости за шаг.
\end{idea}

\subsection{Усложнения транспортной задачи}
\label{subsec:transport-extensions}

Классическая постановка предполагает, что перевозится один однородный продукт,
все маршруты разрешены, а ограничения заданы только запасами и потребностями.
В реальных задачах появляются дополнительные условия.

\paragraph{1. Многопродуктовая модель.}
Перевозится несколько видов продукции. Для продукта $k$ вводятся переменные
$x_{ij}^{(k)}$; ограничения записываются отдельно по каждому виду продукции
либо совместно, с учётом общих ресурсов транспорта:
\[
  \sum_{j} x_{ij}^{(k)} = a_i^{(k)} , \qquad
  \sum_{i} x_{ij}^{(k)} = b_j^{(k)} , \qquad
  \sum_{k} x_{ij}^{(k)} \le d_{ij} ,
\]
где последнее неравенство — общая пропускная способность маршрута.

\paragraph{2. Запрет на поставку.}
Если из пункта $A_i$ нельзя поставлять груз в $B_j$, клетка $(i,j)$
запрещается. Её либо исключают из рассмотрения, либо задают очень большой
тариф $c_{ij} = M$ ($M$ — большое число). Тогда метод минимизации не
использует этот маршрут, пока есть другие допустимые.

\paragraph{3. Обязательная поставка.}
Если по маршруту $(i,j)$ перевозка должна быть, вводится условие
$x_{ij} > 0$; на практике его заменяют ограничением снизу
$x_{ij} \ge K_{ij}$, где $K_{ij} > 0$ — обязательный объём.

\paragraph{4. Поставка не менее заданного количества.}
Ограничение $x_{ij} \ge K_{ij}$ учитывается заранее: обязательную поставку
выделяют и уменьшают соответствующие запас и потребность,
\[
  a_i := a_i - K_{ij} , \qquad b_j := b_j - K_{ij} ,
\]
после чего задача снова становится обычной.

\paragraph{5. Поставка не более заданного количества.}
Ограничение $x_{ij} \le K_{ij}$ означает пропускную способность маршрута.

\paragraph{6. Сменно-суточный характер поставок.}
Если перевозки зависят от времени, вводятся дополнительные индексы:
$x_{ij}^{(t)}$ — перевозка из $A_i$ в $B_j$ в период $t$. Учитываются
ограничения по сменам, суткам, времени доставки и доступности транспорта.

\paragraph{7. Возврат тары и минимизация порожнего пробега.}
Если после доставки нужно вернуть тару или транспорт, добавляются обратные
перевозки $x_{ji}^{\text{back}}$; может ставиться задача минимизации порожнего
пробега, когда транспорт не должен возвращаться пустым.

\begin{idea}
Первые пять усложнений не выводят задачу из класса транспортных: запрет и
верхняя граница — это ограничения на отдельные клетки (их обрабатывают
«запрещённым» тарифом или переходом к транспортной задаче с пропускными
способностями), а нижняя граница снимается сдвигом переменной
$x_{ij}' = x_{ij} - K_{ij}$. Пункты~6 и~7 увеличивают размерность и, как
правило, уже требуют решения общей задачи ЛП симплекс-методом
(\S\ref{subsec:simplex-scheme}), а не табличного метода потенциалов.
\end{idea}
